\scanpage{16}
\mspage{-0.18-}

one $x$. To illustrate this, we give an example and a lemma which suggest opposite conclusions (the example shattered over-optimistic conjectures that the author had entertained!).

\uline{Example} \uline{0.4}: Of a semi-simple group $G$ over $\mathbb C$, the complex numbers, acting on a non-singular quasi-affine scheme of finite type over $\mathbb C$, such that:

i) all stabilizers are reduced to $e$ itself, i.e. the action is set-theoretically \figure{\begin{tikzpicture}[baseline=(f.base)]\node[draw,ellipse,inner xsep=2pt,inner ysep=6pt] (f) {free};\end{tikzpicture}},

ii) a geometric quotient exists, which is also a scheme of finite type over $\mathbb C$; hence (cf. \S4) the action is separated,

iii) the action is not proper, and, in particular, \uline{not: free} (algebro-geometrically).
\margin{\figure{\begin{tikzpicture}[x=1cm,y=1cm]\draw (0,0) .. controls (0.05,.2) and (.12,-.08) .. (.25,.02) .. controls (.4,.12) and (.5,.04) .. (.65,.12);\node at (.9,.35){!};\end{tikzpicture}}}

\uline{Proof}: Take $G=\mathrm{SL}(2)$. Let $V_n$ stand for the $(n+1)$-dimensional affine space whose closed points are homogeneous forms in $X$ and $Y$ of degree $n$ over $\mathbb C$: then $\mathrm{SL}(2)$ acts on $V_n$ by means of substitutions in $X$ and $Y$. Define a reduced subscheme $X\subset V_1\times V_4$ by means of
\[
\begin{aligned}
(F_1,F_4)\in X_{\mathbb C}\Longleftrightarrow\quad &\text{(a) }F_1\ne0\\
&\text{(b) }F_4\text{ is the square of a homogeneous}\\
&\qquad\text{quadratic form of discriminant }1.
\end{aligned}
\]
One checks that $X$ is non-singular, invariant under $\mathrm{SL}(2)$, and that no non-trivial $\alpha\in \mathrm{SL}(2)_{\mathbb C}$ leaves fixed any point of $X$. Define $\phi:X\longrightarrow\mathbb A^1$ by
\mspage{15}
